\section{ظرفیت وابسته به ولتاژ و مدل ساده \lr{MOS}}
\subsection{معنای \lr{C-V} و تغییر بار}
ظرفیت بیان می‌کند که برای تغییر اختلاف ولتاژ دو رسانا، چه مقدار بار باید جابه‌جا
شود. برای خازن خطی با ظرفیت ثابت \(C\)، تغییر بار و اختلاف ولتاژ متناسب‌اند
\begin{equation}
 \Delta Q=C\Delta V.
 \label{eq:cv-linear-charge}
\end{equation}
در افزاره \lr{MOS}، وضعیت کانال و ضخامت ناحیه \lr{Depletion} با بایاس تغییر
می‌کنند. در نتیجه، ظرفیت نیز به ولتاژ وابسته است. نام \lr{C-V Characteristics}
همین وابستگی را بیان می‌کند \cite[آغاز بخش \secnum{۲.۳}]{weste}.

برای ظرفیت غیرخطی، شیب نمودار بار نسبت به ولتاژ در نقطه کار را تعریف می‌کنیم
\begin{equation}
 C(V)=\frac{dQ}{dV},\qquad dQ=C(V)\,dV.
 \label{eq:cv-differential}
\end{equation}
مشتق \(dQ/dV\) یعنی نسبت تغییر بسیار کوچک بار به تغییر بسیار کوچک ولتاژ.
واحد آن \unit{\coulomb\per\volt}، برابر \unit{\farad} است.
برای تغییر کوچک پیرامون \(V_0\)، رابطه
\(\Delta Q\approx C(V_0)\Delta V\) کاربرد دارد.
برای تغییر بزرگ، باید ظرفیت همه بخش‌های مسیر ولتاژ را در نظر گرفت.

اگر ولتاژ از \(V_1\) به \(V_2\) تغییر کند، جمع پیوسته تغییرهای بار می‌دهد
\begin{equation}
 \Delta Q=\int_{V_1}^{V_2}C(V)\,dV,\qquad
 C_{\mathrm{eff},Q}=\frac{\Delta Q}{V_2-V_1}
 =\frac1{V_2-V_1}\int_{V_1}^{V_2}C(V)\,dV.
 \label{eq:cv-charge-equivalent}
\end{equation}
\(C_{\mathrm{eff},Q}\) یک ظرفیت ثابت معادل از نظر بار جابه‌جاشده در همین بازه است.
این استخراج از تعریف ظرفیت است. کتاب برای تخمین دستی نیز از ظرفیت متوسط
طی گذار استفاده می‌کند. تغییر بازه یا حرکت هم‌زمان پایانه دیگر می‌تواند مقدار
معادل را تغییر دهد.

برای مثال فرضی \(C(V)=C_a(1-aV)\)، \(C_a=\qty{2}{\femto\farad}\)،
\(a=\qty{0.2}{\per\volt}\) و بازه صفر تا \qty{1}{\volt}، ظرفیت در ابتدا
\qty{2}{\femto\farad} و در انتها \qty{1.6}{\femto\farad} است.
با جمع مساحت زیر خط، میانگین \qty{1.8}{\femto\farad} و تغییر بار
\qty{1.8}{\femto\coulomb} می‌شود. چون در این مثال \(C(V)\) خطی است، میانگین
دو مقدار انتهایی با میانگین انتگرالی برابر است.
در رابطه غیرخطی، میانگین دو سر بازه یک تقریب است.

\subsection{مدار معادل پنج ظرفیت و پایانه مرجع}
کتاب در شکل \secnum{۲.۱۳} پنج ظرفیت عمده را نشان می‌دهد. سه سهم
\(C_{GS},C_{GD},C_{GB}\) به \lr{Gate}، و دو سهم \(C_{SB},C_{DB}\)
به پیوندهای \lr{Diffusion} و \lr{Body} مربوط‌اند.
شکل \ref{fig:cv-network} این ارتباط را بدون جزئیات ترانزیستور بازترسیم می‌کند.
این نمایش، مدل خازنی آموزشی بایاس‌شده است \cite[صفحه ۷۳]{weste}.
\begin{figure}[ht]
\centering
\begin{latin}
\begin{circuitikz}[font=\small]
\draw (0,3) node[circ] (G) {} node[above] {Gate};
\draw (-3,0) node[circ] (S) {} node[left] {Source};
\draw (3,0) node[circ] (D) {} node[right] {Drain};
\draw (0,-3) node[circ] (B) {} node[below] {Body};
\draw (G) to[C,l=$C_{GS}$] (S);
\draw (G) to[C,l_=$C_{GD}$] (D);
\draw (S) to[C,l=$C_{SB}$] (B);
\draw (D) to[C,l_=$C_{DB}$] (B);
\draw (G) to[C,l=$C_{GB}$] (B);
\end{circuitikz}
\end{latin}
\caption{ارتباط پایانه‌ای ظرفیت‌های اصلی در مدل آموزشی کتاب.
سهم‌های \lr{Intrinsic} و \lr{Overlap} در \(C_{GS}\) و \(C_{GD}\) جمع شده‌اند
\cite[شکل \secnum{۲.۱۳}]{weste}.}
\label{fig:cv-network}
\end{figure}

برای یک خازن ثابت میان گره \(A\) و گره \(B\)، بار سمت \(A\) برابر
\(Q_A=C(V_A-V_B)\) و جریان ورودی آن سمت برابر است با
\begin{equation}
 i_A=C\left(\frac{dV_A}{dt}-\frac{dV_B}{dt}\right).
 \label{eq:cv-two-terminal-current}
\end{equation}
اگر \(B\) ثابت بماند، \(dV_B/dt=0\) و ظرفیت مانند بار متصل به مرجع ثابت دیده
می‌شود. اتصال \lr{Body} به تغذیه ثابت نیز برای تحلیل تغییرات ولتاژ چنین نقشی دارد.
اگر هر دو گره حرکت کنند، اختلاف تغییرهای آن‌ها در جریان مؤثر است.

ظرفیت‌های متصل به یک گره را هنگامی مستقیم جمع می‌کنیم که سمت دیگرشان در تحلیل
مورد نظر مرجع ثابت باشد یا اثر حرکت آن جداگانه منظور شده باشد.
برای نمونه، ظرفیت \lr{Drain} به \lr{Body} ثابت، سهمی از بار خروجی است.
ظرفیت \lr{Gate} مرحله بعد، از دید خروجی مرحله قبل بار ورودی آن مرحله محسوب می‌شود.

\subsection{مدل صفحه‌ای \lr{Gate} و ظرفیت بر واحد عرض}
در مدل ساده بخش \secnum{۲.۳.۱}، ظرفیت \lr{Gate} روی کانال از رابطه
\begin{equation}
 C_g\approx C_0=WLC_{\mathrm{ox}},\qquad
 C_{\mathrm{ox}}=\frac{\varepsilon_{\mathrm{ox}}}{t_{\mathrm{ox}}}
 \label{eq:cv-simple-gate}
\end{equation}
تخمین زده می‌شود. \(C_{\mathrm{ox}}\) ظرفیت بر واحد سطح و \(C_g\) ظرفیت کل است.
کانال پایانه مستقل بیرونی ندارد. کتاب برای تخمین مقدماتی، صفحه پایین خازن را
به \lr{Source} نسبت می‌دهد و از \(C_{GS}\) استفاده می‌کند. در مدل تفصیلی،
این نسبت‌دادن به ناحیه کاری وابسته خواهد بود
\cite[رابطه \secnum{۲.۱۲} و اسلاید ۲۴]{weste,mos}.

در فناوری با طول ثابت، می‌توان \(L\) را در ضریب ظرفیت بر واحد عرض جذب کرد
\begin{equation}
 c_w=C_{\mathrm{ox}}L
 =\frac{\varepsilon_{\mathrm{ox}}L}{t_{\mathrm{ox}}},\qquad C_g\approx c_wW.
 \label{eq:cv-width-coefficient}
\end{equation}
نماد \(c_w\) همان \(C_{\mathrm{permicron}}\) در کتاب است، با واحد رایج
\unit{\femto\farad\per\micro\metre}.
اگر \(c_w=\qty{2}{\femto\farad\per\micro\metre}\) و
\(W=\qty{0.6}{\micro\metre}\) باشند، ظرفیت \(C_g\approx\qty{1.2}{\femto\farad}\)
می‌شود. واحد طول در ضرب حذف می‌شود.

از نظر تبدیل واحد، \(\qty{1}{\femto\farad\per\micro\metre}=10^{-9}\,
\unit{\farad\per\metre}\) است. برای نمونه فرضی،
\(C_{\mathrm{ox}}=\qty{4}{\femto\farad\per\micro\metre\squared}\) و
\(L=\qty{0.25}{\micro\metre}\) ضریب \(c_w=\qty{1}{\femto\farad\per\micro\metre}\)
می‌دهند. با \(W=\qty{2}{\micro\metre}\)، ظرفیت \qty{2}{\femto\farad} است.

اگر \(L\) و \(t_{\mathrm{ox}}\) با ضریب یکسان کوچک شوند و گذردهی ثابت بماند،
نسبت \(L/t_{\mathrm{ox}}\) و در نتیجه \(c_w\) ثابت می‌ماند.
اگر عرض نیز با همان ضریب کوچک شود، ظرفیت کل به نسبت عرض کاهش می‌یابد.
کتاب مقدار حدود \qty{2}{\femto\farad\per\micro\metre} را برای فرآیندهای
قدیمی‌تر و حدود \qty{1}{\femto\farad\per\micro\metre} را در نمونه‌های
\qty{90}{\nano\metre} و \qty{65}{\nano\metre} ذکر می‌کند.
این اعداد تخمین‌های تاریخی همان مرجع هستند \cite[روابط \secnum{۲.۱۳} و \secnum{۲.۱۴}]{weste}.

\subsection{مدل ساده \lr{Diffusion} و ارتباط با بار خروجی}
پیوند \lr{p-n} میان \lr{Source} یا \lr{Drain} و \lr{Body} در بایاس معکوس
ناحیه \lr{Depletion} دارد. دو ناحیه رسانا در دو سوی این ناحیه کم‌حامل، رفتار
خازنی ایجاد می‌کنند. پهنای ناحیه با ولتاژ تغییر می‌کند، بنابراین ظرفیت به بایاس
وابسته است \cite[اسلاید ۲۳ و صفحه ۶۹]{mos,weste}.
اصطلاح \lr{Diffusion Capacitance} در این فصل نام رایج طراحان مدار برای همین
ظرفیت پیوندی \lr{Depletion} است. این کاربرد در پانوشت صفحه ۶۹ توضیح داده شده است.

کتاب برای تخمین دستی، \(C_{SB}\) و \(C_{DB}\) ناحیه‌های دارای \lr{Contact}
را هم‌مرتبه ظرفیت \lr{Gate} می‌گیرد و در بسیاری از مثال‌ها ضریب حدود
\qty{1}{\femto\farad\per\micro\metre} به‌ازای عرض را به کار می‌برد.
ضریب مناسب هر فناوری و هندسه باید از داده همان مدل برداشته شود.
هندسه ناحیه مشترک و بدون \lr{Contact} در تحلیل دقیق اثر دارد
\cite[صفحه ۷۰]{weste}.

در مثال فرضی، یک \lr{Inverter} خروجیِ متصل به دو \lr{Inverter} مشابه بعدی دارد.
عرض‌ها \(W_n=\qty{1}{\micro\metre}\)، \(W_p=\qty{2}{\micro\metre}\) و
تمام ضریب‌های ساده \lr{Gate} و \lr{Drain} برابر
\qty{1}{\femto\farad\per\micro\metre} فرض می‌شوند. ظرفیت سیم
\qty{0.5}{\femto\farad} است. هر ورودی بعدی \qty{3}{\femto\farad} و مجموع
\lr{Drain}های خروجی راه‌انداز نیز \qty{3}{\femto\farad} می‌دهد. پس
\begin{equation}
 C_L\approx2(1+2)+(1+2)+0.5=\qty{9.5}{\femto\farad}.
 \label{eq:cv-load-example}
\end{equation}
اعداد داخل جمع، ظرفیت در واحد \unit{\femto\farad} هستند.
\lr{Source}های راه‌انداز به ریل ثابت متصل‌اند و در این مدل، سهم پیوندی مستقیم
خروجی از \lr{Drain}های آن می‌آید. با ولتاژ تغییر \qty{1}{\volt} و جریان متوسط
فرضی \qty{100}{\micro\ampere}، تخمین \(C_L\Delta V/I\) برابر
\qty{95}{\pico\second} می‌شود. این تخمین از ضریب‌های ساده و مرجع‌های ثابت استفاده می‌کند.

\subsection{تمرین‌های پاسخ‌دار مدل ساده}
\textbf{تمرین ۱.} اگر \(C_{\mathrm{ox}}=\qty{5}{\femto\farad\per\micro\metre\squared}\)،
\(L=\qty{0.2}{\micro\metre}\) و \(W=\qty{0.8}{\micro\metre}\) باشند،
ظرفیت چقدر است؟ ابتدا ضریب بر واحد عرض و سپس ظرفیت کل را می‌یابیم.
\[
 \begin{aligned}
 c_w&=5(0.2)=\qty{1}{\femto\farad\per\micro\metre},\\
 C_g&=1(0.8)=\qty{0.8}{\femto\farad}.
 \end{aligned}
\]

\textbf{تمرین ۲.} در مقیاس‌بندی ایدئال، \(L\)، \(W\) و ضخامت اکسید هر یک
نصف شوند. نسبت ظرفیت جدید به قبلی چیست؟ \(C_{\mathrm{ox}}\) دو برابر و
\(WL\) یک‌چهارم می‌شود، پس \(C_g\) نصف خواهد شد.
این نتیجه به فرض ثابت‌بودن گذردهی و مدل صفحه‌ای وابسته است.

\textbf{تمرین ۳.} سه ظرفیت \qty{2}{\femto\farad}، \qty{3}{\femto\farad} و
\qty{1}{\femto\farad} از یک گره به مرجع ثابت وصل‌اند. برای تغییر
\qty{0.8}{\volt}، بار جابه‌جاشده چقدر است؟ ظرفیت کل
\(2+3+1=\qty{6}{\femto\farad}\) و تغییر بار
\(6(0.8)=\qty{4.8}{\femto\coulomb}\) می‌شود.
\clearpage
